% D7 草稿 v2.0（2026-11-01）: 按新版四条 RQ(docs/RQ_定义与映射_v1.0.md, commit 4579859)对齐。
% 相对 v1.0(D8 对账版)的改动:
%   [1] Findings 由 3 条扩为 4 条, 逐条挂 RQ 标签: ①RQ1 赢家翻转(并入官方 5:1 代价维度)
%       ②RQ1/RQ2 买到召回、精度代价随数据集 ③RQ2 accuracy 掩盖一切(补预注册规则 10 单元命中)
%       ④RQ4 新样本决策(新增, 数字= holdout 一次性评估, D9 已核对)。
%   [2] Limitations 增第四条: RQ4 holdout 单次评估、与训练同分布、上线需漂移监控。
%   [3] 给 E 的验证摘要注释块由 3 句扩为 4 句(新增 RQ4 一句), 供 Abstract 使用。
%   [4] 加 \label{sec:conclusion}(sec6 v2.0 的 RQ4 小节引用它)。
%   其余正文与 v1.0 逐字一致(D8 已全量对账)。
% 写回方法(待用户核验后): cp "docs/D7_sec7_conclusion_draft_v2.0.tex" overleaf/sections/sec7_conclusion.tex (去掉本注释头)
% 源提交: 053de96(D7初稿)+af8d4e3(Holm勘误)+6944ab5(D8勘误E1)+4579859(新RQ决议)

\section{Conclusion and Future Work}
\label{sec:conclusion}
% D drafts this file; E helps align conclusions with measured results and limitations.
%
% === 4-sentence verified results summary for E (abstract / sec1 use; 2026-11-01, D, RQ-aligned v2.0) ===
% 1. (RQ1) On credit-g (30% positives) the winner is cost-sensitive logistic regression trained with
%    the official 5:1 cost (cv10 F1 0.597, recall 0.863 vs 0.470 plain; pooled cost 0.884 -> 0.514),
%    while on sick (6.1% positives) it is SMOTE-resampled random forest (F1 0.897, recall 0.888)
%    --- the winner flips across datasets and none of size, attribute composition, missingness, or
%    imbalance alone predicts it.
% 2. (RQ1/RQ2) All three imbalance strategies reliably buy recall (9 of 12 pre-specified paired
%    comparisons significant, Nadeau--Bengio corrected), but F1 gains flip sign on sick for
%    SMOTE x (DT/LR): the precision price, not the recall gain, is dataset-dependent.
% 3. (RQ2) Accuracy hides everything this project measures: ZeroR reaches 70%/94% accuracy with
%    zero recall, and the pre-registered trap rule (acc >= 0.70, recall <= 0.40) flags ten archived
%    units; decision metrics must come from the task asymmetry (official cost on credit-g,
%    positive recall on sick), with AUC/PR as ranking evidence.
% 4. (RQ4) On the locked hold-outs, evaluated once, the champions make quantified, defensible
%    decisions on unseen samples: lr_cost catches 82% of bad applicants at cost 0.57 (best plain
%    0.89, always-approve 1.50) and rf_smote catches 66/69 patients (recall 0.957, precision 0.868).
% === end summary ===

\subsection{Findings}

This project compared four classifiers (DT, RF, LR, SVC) under three imbalance-handling strategies (SMOTE, cost-sensitive learning, threshold adjustment) on two locked UCI benchmarks: the moderately imbalanced, fully observed credit-g (30\% positives) and the extremely imbalanced, noisy sick (6.1\% positives, 29.9\% of rows with missing values). Under a frozen 10-fold cross-validation protocol with overlap-corrected paired tests, four findings --- one per research question --- are robust:

\begin{enumerate}
  \item \textbf{(RQ1) The winner flips across datasets.} credit-g rewards \emph{linear + cost-sensitive} learning (\texttt{lr\_cost}, F1 $0.597$, with the official 5:1 cost folded into training); sick rewards \emph{ensemble + resampling} (\texttt{rf\_smote}, F1 $0.897$). Four dataset properties --- size, attribute composition, missingness, and imbalance --- jointly explain the flip; the single-variable reading ``more imbalance needs more resampling'' predicts the opposite of what we observe (Section~\ref{sec:results-flip}).
  \item \textbf{(RQ1/RQ2) Strategies buy recall; datasets set the precision price.} Recall improves in all twelve pre-specified strategy-vs-plain comparisons (nine significant), but F1 flips sign on sick precisely where SMOTE meets linear or small-tree models ($-0.166$ for LR, $p=0.002$). Imbalance interventions are not free and not portable.
  \item \textbf{(RQ2) Accuracy hides everything this project measures.} ZeroR reaches 70\%/94\% accuracy on credit-g/sick with zero recall; the pre-registered trap rule (accuracy $\geq 0.70$, recall $\leq 0.40$) flags ten archived units, and plain tree models sit within a few points of ZeroR accuracy while detecting a quarter of the minority class. Evaluating minority-class pipelines by accuracy alone would have ranked every serious model below the trivial one; decision metrics must instead come from the task's asymmetry --- the official cost on credit-g, positive-class recall on sick --- with ROC/PR curves as the ranking evidence.
  \item \textbf{(RQ4) The learnt models make reliable, quantified decisions on unseen samples.} Evaluated once on the locked hold-outs, the credit-g champion catches 82\% of bad applicants and cuts the official cost to 0.57 (best plain 0.89, always-approve 1.50), and the sick champion catches 66 of 69 patients (recall 0.957) at precision 0.87. Both decision regimes --- loss control and miss prevention --- survive contact with genuinely new data with error profiles a deployment decision can be argued from (Section~\ref{sec:results-holdout}).
\end{enumerate}

\subsection{Limitations}

Four limitations bound these conclusions. First, the \textbf{ensemble study is exploratory}: it was added after the main matrix was frozen, is kept in a separate results namespace, and only two comparisons were pre-specified; its cells are excluded from the main statistical family and from Holm correction, and the sick champion-vs-ensemble gap ($0.937$ vs $0.897$) is descriptive, untested. Second, \textbf{SMOTE operates on one-hot-encoded categorical attributes}, where linear interpolation creates fractional categories that are not realisable data points; this is a plausible mechanism for the observed SMOTE$\times$LR collapse but was not manipulated experimentally. Third, the design uses \textbf{one locked split seed (42)} with deliberately conservative overlap-corrected tests; effects reported as non-significant (e.g.\ credit-g strategy gains, $p=0.03$--$0.07$ raw) should be read as ``not established at our conservative threshold'', not as evidence of absence. Fourth, the \textbf{RQ4 evidence is a single evaluation per dataset on same-period hold-outs} (300 and 1{,}132 instances): it certifies decision quality under the training distribution only --- drift between training and deployment is outside what a one-shot hold-out can measure.

\subsection{Future Work}

The clearest next step is \textbf{SMOTE-NC} (or another mixed-type synthesiser), which interpolates only in the numeric subspace and samples categorical values by mode --- directly targeting the interpolation-noise mechanism identified above; the prediction it tests is that LR/DT SMOTE losses on sick shrink while RF is unaffected. Beyond that: threshold selection by expected-cost minimisation on the credit-g cost matrix (rather than F1 on development folds); calibrated probabilities to make threshold and cost interventions comparable across models; a nested-CV replication over multiple split seeds to attach uncertainty to the winner-flip itself; and, operationally, a deployment loop for the RQ4 champions --- drift monitoring with scheduled re-validation, since the hold-out certifies only the collection-period distribution.
